\documentclass[10pt,twocolumn,letterpaper]{article}

% --- Packages ---
\usepackage[utf8]{inputenc}
\usepackage[T1]{fontenc}
\usepackage{lmodern}
\usepackage[margin=0.65in,top=0.75in,bottom=0.65in]{geometry}
\usepackage{amsmath}
\usepackage{amssymb}
\usepackage{graphicx}
\usepackage{booktabs}
\usepackage{hyperref}
\usepackage{xcolor}
\usepackage{caption}
\usepackage{natbib}
\usepackage{microtype}

% Compact spacing
\setlength{\parskip}{1.5pt plus 1pt}
\setlength{\parindent}{0.5em}
\setlength{\abovecaptionskip}{4pt}
\setlength{\belowcaptionskip}{2pt}
\setlength{\textfloatsep}{6pt}
\setlength{\floatsep}{4pt}
\setlength{\dbltextfloatsep}{6pt}

% Compact sections
\makeatletter
\renewcommand\section{\@startsection{section}{1}{\z@}{5pt}{2pt}{\normalfont\bfseries\normalsize}}
\renewcommand\subsection{\@startsection{subsection}{2}{\z@}{3pt}{1.5pt}{\normalfont\bfseries\small}}
\makeatother

% Compact abstract
\renewenvironment{abstract}{\small\noindent\textbf{Abstract.}\enspace\ignorespaces}{\par\vspace{4pt}}

\hypersetup{colorlinks=true, linkcolor=blue!70!black, citecolor=blue!70!black, urlcolor=blue!70!black}

\title{\vspace{-8pt}\textbf{Scrutinizer: A Biologically-Grounded Foveated Vision Simulator for Web Interfaces}\vspace{-4pt}}

\author{
  Andy Edmonds\\[-2pt]
  {\small Independent Researcher, San Francisco, CA}\\[-2pt]
  {\small\texttt{andyed@alum.cmu.edu}}
}

\date{\vspace{-12pt}}

\begin{document}
\maketitle
\vspace{-6pt}

% ============================================================
\begin{abstract}
We present Scrutinizer, an open-source, real-time gaze-contingent foveated rendering system for live web content. A 1157-line GLSL shader implements a pipeline mapped to the visual pathway (retinal ganglion cells~$\rightarrow$ LGN~$\rightarrow$ V1~$\rightarrow$ V4), with four swappable JavaScript modules aligned to biological subsystems. The key technical contribution is a Difference-of-Gaussians (DoG) band decomposition exploiting the hardware MIP chain for frequency-selective peripheral filtering---high-frequency detail attenuates first while low-frequency structure persists, mirroring cortical magnification~\citep{rovamo1979}. A reference implementation of FOVI~\citep{blauch2026} enables direct comparison of cortical magnification strategies within the same framework. Open source (MIT), 60\,fps, 138 unit tests.
\end{abstract}

% ============================================================
\section{Introduction}

Peripheral vision is not blurry---it is \emph{differently encoded}. The retina compresses ${\sim}10^7$ photoreceptor signals into ${\sim}10^6$ optic nerve fibers, selectively preserving low-frequency structure while filtering high-frequency detail~\citep{curcio1990, dacey1993}. This is compute demand management: limited bandwidth allocated to information most likely to guide the next saccade.

Existing gaze-contingent tools---BubbleView~\citep{kim2017}, ViewSer~\citep{lagun2011}, ScreenMasker---apply Gaussian blur as a peripheral proxy, uniformly destroying spatial structure and preventing layout-level judgments that real peripheral vision supports. These tools also operate on static images, not live interactive content.

\begin{table}[t]
\centering
\caption{Comparison with existing gaze-contingent tools.}\vspace{-2pt}
\label{tab:comparison}
\footnotesize
\begin{tabular}{@{}lccccc@{}}
\toprule
\textbf{Tool} & \rotatebox{60}{\textbf{Live}} & \rotatebox{60}{\textbf{Bio Model}} & \rotatebox{60}{\textbf{Crowding}} & \rotatebox{60}{\textbf{Saliency}} & \rotatebox{60}{\textbf{Open}} \\
\midrule
BubbleView~\citeyearpar{kim2017} & -- & Gauss. & -- & -- & \checkmark \\
ViewSer~\citeyearpar{lagun2011} & SERP & Gauss. & -- & -- & -- \\
ScreenMasker & -- & Gauss. & -- & -- & \checkmark \\
FOVI~\citeyearpar{blauch2026} & \checkmark & CMF & -- & -- & \checkmark \\
\textbf{Scrutinizer} & \checkmark & \textbf{DoG} & \checkmark & \checkmark & \checkmark \\
\bottomrule
\end{tabular}
\end{table}

Scrutinizer contributes \textbf{(1)} a modular architecture where each module maps to a biological subsystem (Table~\ref{tab:architecture}), enabling independent swapping---replace the mouse-based gaze proxy with an eye tracker by changing one module---and \textbf{(2)} a swappable pipeline enabling side-by-side comparison of peripheral filtering strategies, including a reference implementation of FOVI cortical magnification~\citep{blauch2026}.

% ============================================================
\section{Architecture: Biology $\rightarrow$ Computation}

The pipeline is a single-pass fragment shader (1157~LOC, GLSL~ES~3.0) with 27 uniforms organized into biological stages, fed by four JavaScript modules (Table~\ref{tab:architecture}).

\begin{table}[t]
\centering
\caption{Module architecture mapped to biological subsystems.}\vspace{-2pt}
\label{tab:architecture}
\footnotesize
\begin{tabular}{@{}p{1.4cm}p{1.7cm}p{3.5cm}@{}}
\toprule
\textbf{Biology} & \textbf{Module (LOC)} & \textbf{Function} \\
\midrule
Oculomotor & GazeModel (166) & Velocity, fixation, saccadic suppression \\
Retinal GC & ContentAnalysis (361) & DoG center-surround via MIP chain \\
LGN gating & ContentAnalysis & Structure mask + saliency gate \\
V1 features & Orchestrator (536) & Noise crowding, grid scramble \\
V4 color & Renderer (548) & Oklab desat., chrom.\ aberration \\
Saliency & Worker (489) & Oklab DoG + face det.\ + structure \\
\bottomrule
\end{tabular}
\end{table}

\subsection{DoG Band Decomposition}

Retinal ganglion cells have center-surround receptive fields modeled as DoG filters~\citep{kuffler1953}, with field size growing with eccentricity (M-scaling). We decompose the hardware MIP chain into a Laplacian pyramid:
\vspace{-2pt}
\begin{equation}
\text{band}_k = \text{textureLod}(\mathit{tex}, \mathit{uv}, k) - \text{textureLod}(\mathit{tex}, \mathit{uv}, k{+}1)
\end{equation}
\vspace{-4pt}

Four bands span 1--2\,px (serifs) to 8--16\,px (buttons), each weighted by a smoothstep rolloff with cutoffs following a geometric progression ($0.3, 0.6, 1.2, 2.4 \times E_2$), approximating cortical magnification~\citep{rovamo1979}. The DC residual always passes through. Result: peripheral text shows letter shapes (low-frequency visible, identity unreadable) rather than uniform fog. Cost: 5 \texttt{textureLod} calls vs.\ 1, but near-zero practical overhead---higher MIP levels are cheaper, already cached, and arithmetic is trivial.

\subsection{Saliency-Driven Bandwidth Allocation}

The LGN's synaptic inputs are ${\sim}90\%$ modulatory, with ${\sim}30\%$ from V1 feedback~\citep{sherman2002}; spatial attention enhances LGN gain~\citep{mcalonan2008}. We model this as saliency gating: a Web Worker computes a saliency map (Oklab opponent-channel DoG + face detection + Gestalt structure masks), uploaded as a GPU texture. High-saliency regions receive up to 70\% of rendering bandwidth.

\subsection{V1 Crowding and V4 Color}

Crowding is feature binding failure, not blur~\citep{pelli2008}. We model it via simplex noise domain warping and grid-based feature scrambling~\citep{rosenholtz2012}. V4 applies Oklab desaturation~\citep{ottosson2020} for rod dominance and chromatic aberration for M/P asynchrony~\citep{livingstone1988}.

\subsection{Extensibility Branch Points}

The architecture exposes five branch points where researchers substitute implementations without modifying the core pipeline (Table~\ref{tab:branch}).

\begin{table}[h]
\vspace{-2pt}
\centering
\caption{Extensibility branch points.}\vspace{-2pt}
\label{tab:branch}
\footnotesize
\begin{tabular}{@{}p{1.3cm}p{2.1cm}p{3.2cm}@{}}
\toprule
\textbf{Branch} & \textbf{Interface} & \textbf{Swap examples} \\
\midrule
Gaze input & \texttt{update(x,y)} $\rightarrow$ pos, vel, fixation & Tobii SDK, WebGazer.js, scanpath replay \\
Saliency & bitmap $\rightarrow$ heatmap texture & Itti-Koch, DeepGaze~II, ONNX model \\
Crowding & integer uniform selects V1 type & Portilla-Simoncelli, oriented DoG \\
Structure & DOM blocks $\rightarrow$ RGBA texture & CV segmentation, semantic embed \\
Pipeline & JSON registry & Zero-code new profiles \\
\bottomrule
\end{tabular}
\end{table}

\subsection{Cortical Magnification: FOVI Comparison}

\citet{blauch2026} propose FOVI, a cortical magnification transform for vision models: $\mathit{mipLevel}(e) = \log_2(1 + e/a)$, $a{=}2.78^{\circ}$. We implemented FOVI alongside Scrutinizer's DoG pipeline, enabling direct comparison. At matched eccentricities, Scrutinizer's effective cutoff falls 3.5--4.6$\times$ lower than FOVI's---the gap represents crowding, receptive-field growth, and contrast sensitivity loss that Scrutinizer bakes into its shader stages (V1 noise, saliency gating, Oklab desaturation) while FOVI leaves to the downstream neural network.

Both systems share an open question: the desaturation transfer function. Scrutinizer uses smoothstep in Oklab space; FOVI applies exponential decay in RGB. Neither approach is definitively grounded---\citet{mullen2002} show red-green sensitivity declining faster than blue-yellow with eccentricity, suggesting a dual-sigma model per opponent channel rather than a uniform curve.

% ============================================================
\section{Development Process}

The system was developed through iterative AI-human collaboration (Claude, Anthropic) where the agent served as \emph{research collaborator} maintaining biological grounding.

\textbf{Cross-process bug detection.}
The saliency worker computed Oklab opponent channels on BGRA-ordered pixels from Electron's \texttt{capturePage()}. The agent traced the byte-order mismatch across 4 files and 3 processes (main $\rightarrow$ renderer $\rightarrow$ Worker)---a class of bug where human code review is unlikely to catch the cross-boundary semantic error.

\textbf{Conceptual reframing.}
Documentation described ``degradation with preservation.'' The human identified the biological analog is \emph{selective bandwidth allocation}; the agent then swept 12 files (${\sim}$30 edits) to align language, illustrating the agent as conceptual partner rather than code generator.

The modular architecture \emph{enables} this workflow: the agent reasons per biological module, not across a monolithic shader.

% ============================================================
\section{Status \& Future Work}

\textbf{Implemented:}
DoG band decomposition; saliency bandwidth allocation; noise crowding and grid scrambling; Oklab desaturation and chromatic aberration; foveal calibration; mobile emulation; FOVI cortical magnification reference implementation; hot-switchable pipeline profiles; 138 tests; golden capture regression.

\textbf{Future:}
Dual-sigma color decay per opponent channel~\citep{mullen2002}; oriented DoG bands (oblique effect~\citep{appelle1972}); calibrated visual angles; eye tracker integration (GazeModel branch point); Portilla-Simoncelli mongrel synthesis (V1 branch point).

\section{Availability}

Open source (MIT), Electron~30.5, macOS signed. Architecture docs, literature review, and a 14-project grad student backlog ship with the repository.

% ============================================================
\vspace{2pt}
\noindent\textbf{Acknowledgments.}
{\small System architecture, documentation, and this paper were developed through collaboration with Claude (Anthropic, Opus~4.6) via Claude Code. The AI agent exclusively authored the grad student project backlog~\citep*{hcibib}. AE directed architecture, identified conceptual framings, and validated perceptual accuracy.}

\vspace{2pt}
\begingroup
\footnotesize
\setlength{\bibsep}{1pt}
\bibliographystyle{plainnat}
\bibliography{references}
\endgroup

\end{document}
